\begin{abstract}
The proliferation of high-fidelity GAN-generated faces has made deepfake
imagery increasingly difficult to detect visually.
We investigate whether combining spatial (RGB) and spectral (FFT) cues in a
unified deep neural network outperforms either modality alone.
Three strategies are benchmarked under identical conditions:
(\emph{i})~an RGB-only ResNet18 fine-tuned from ImageNet weights;
(\emph{ii})~an FFT-only ResNet18 trained from scratch on log-magnitude
frequency spectra;
(\emph{iii})~a Hybrid dual-branch model that late-fuses features from both
domains.
All experiments use a balanced 30{,}000-image subset of the publicly
available 140K Real and Fake Faces dataset (StyleGAN2 vs.\ FFHQ/Flickr).
On a 4{,}500-image held-out test set, the Hybrid model achieves
\textbf{99.5\% accuracy} and \textbf{AUC\,=\,0.9998}, surpassing the
strong RGB baseline (99.0\%, AUC\,=\,0.9975) and the FFT-only model
(61.0\%, AUC\,=\,0.6412).
Crucially, the Hybrid attains perfect recall (zero false negatives).
Gradient-weighted Class Activation Mapping (Grad-CAM) confirms that the RGB
branch attends to semantically meaningful facial regions.
Code, metrics, and figures are publicly available on GitHub.
\end{abstract}
